\documentclass[11pt,a4paper]{ctexart}

\usepackage[a4paper,margin=2.4cm]{geometry}
\usepackage{amsmath,amssymb,mathtools,bm}
\usepackage{booktabs,array,longtable}
\usepackage{xcolor}
\usepackage{listings}
\usepackage{hyperref}
\usepackage{enumitem}

\hypersetup{colorlinks=true,linkcolor=blue!50!black,urlcolor=blue!60!black}
\setlength{\parindent}{2em}
\setlength{\parskip}{0.25em}
\numberwithin{equation}{section}

\definecolor{codebg}{RGB}{247,248,250}
\definecolor{codegreen}{RGB}{20,110,70}
\definecolor{codeblue}{RGB}{30,70,160}
\lstset{
  language=Matlab,
  basicstyle=\small\ttfamily,
  keywordstyle=\color{codeblue}\bfseries,
  commentstyle=\color{codegreen},
  backgroundcolor=\color{codebg},
  frame=single,
  rulecolor=\color{black!20},
  numbers=left,
  numberstyle=\tiny\color{black!55},
  stepnumber=1,
  numbersep=8pt,
  breaklines=true,
  columns=fullflexible,
  keepspaces=true,
  showstringspaces=false,
  tabsize=4
}

\newcommand{\R}{\mathbb{R}}
\newcommand{\vect}[1]{\bm{#1}}
\newcommand{\dd}{\,\mathrm{d}}
\newcommand{\ii}{\mathrm{i}}

\title{从零理解 Chebyshev 谱配置法\\
\large Lagrange 插值、谱微分矩阵与一维 Poisson 边值问题}
\author{学习手册}
\date{\today}

\begin{document}
\maketitle
\tableofcontents
\newpage

\section{学习目标与总体路线}

本手册从一个具体问题出发：只知道函数在有限个节点上的值，怎样计算这些节点上的导数，并进一步求解微分方程？最终目标是理解并实现
\begin{equation}
  \vect{u}'\approx D\vect{u},
  \qquad
  \vect{u}''\approx D^2\vect{u}.
\end{equation}
这里的关键对象是谱微分矩阵 $D$。整条逻辑链为
\begin{equation*}
\text{节点值}
\longrightarrow \text{全局插值多项式}
\longrightarrow \text{对插值多项式求导}
\longrightarrow \text{微分矩阵}
\longrightarrow \text{代数方程组}.
\end{equation*}

读完后应能回答以下问题：
\begin{enumerate}[label=\arabic*.]
  \item 为什么选择 Lagrange 基函数？
  \item $D_{ij}=\ell_j'(x_i)$ 的行、列分别表示什么？
  \item Chebyshev 节点为什么会产生 $(-1)^{i+j}$？
  \item 为什么端点修正系数为 $c_0=c_N=2$？
  \item 为什么对角元可以取非对角元的负行和？
  \item 为什么二阶导数矩阵是 $D^2$ 而不是 $D.^2$？
\end{enumerate}

\section{从连续函数到节点值}

设函数 $u(x)$ 定义在 $[-1,1]$ 上。计算机不能保存连续函数的全部信息，因此选取 $N+1$ 个互异节点
\begin{equation}
  x_0,x_1,\ldots,x_N
\end{equation}
并保存节点值
\begin{equation}
  u_j:=u(x_j),\qquad j=0,1,\ldots,N.
\end{equation}
组成向量
\begin{equation}
  \vect{u}
  =\begin{bmatrix}u_0&u_1&\cdots&u_N\end{bmatrix}^{\mathsf T}.
\end{equation}

谱配置法不直接对离散数据做局部差分，而是先构造次数不超过 $N$ 的全局多项式 $p_N(x)$，使得
\begin{equation}
  p_N(x_j)=u_j,\qquad j=0,1,\ldots,N.
\end{equation}
然后使用
\begin{equation}
  u'(x_i)\approx p_N'(x_i).
\end{equation}
需要明确：我们对 $p_N$ 的求导是精确的；近似发生在用 $p_N$ 代替原函数 $u$ 这一步。

\section{Lagrange 基函数：让系数直接等于节点值}

\subsection{定义}

第 $j$ 个 Lagrange 基函数定义为
\begin{equation}
  \boxed{
  \ell_j(x)=
  \prod_{\substack{k=0\\k\ne j}}^N
  \frac{x-x_k}{x_j-x_k}}
  \label{eq:lagrange-basis}
\end{equation}
其中 $j$ 固定，$k$ 是连乘指标。

当 $x=x_i$ 且 $i\ne j$ 时，分子中包含因子 $x_i-x_i=0$，因此
\begin{equation}
  \ell_j(x_i)=0.
\end{equation}
当 $x=x_j$ 时，式~\eqref{eq:lagrange-basis} 中每个分子都等于对应分母，因此
\begin{equation}
  \ell_j(x_j)=1.
\end{equation}
统一写成
\begin{equation}
  \boxed{\ell_j(x_i)=\delta_{ij}},
  \qquad
  \delta_{ij}=\begin{cases}1,&i=j,\\0,&i\ne j.\end{cases}
\end{equation}

这个等式只描述有限个节点处的取值，并不表示 $\ell_j(x)$ 恒为零。固定 $j$ 后，$\ell_j$ 在自己的节点处为 $1$，在另外 $N$ 个节点处为 $0$，在节点之间一般不为零。

\subsection{插值多项式}

定义
\begin{equation}
  \boxed{p_N(x)=\sum_{j=0}^N u_j\ell_j(x)}.
  \label{eq:interpolant}
\end{equation}
在节点 $x_i$ 处，
\begin{equation}
  p_N(x_i)=\sum_{j=0}^N u_j\delta_{ij}=u_i.
\end{equation}
因此，式~\eqref{eq:interpolant} 自动穿过全部节点。选择 Lagrange 基的思想动机正是：展开系数 $u_j$ 直接就是具有物理意义的节点值，不需要再从单项式系数反解节点值。

\section{微分矩阵的定义与意义}

对式~\eqref{eq:interpolant} 求导：
\begin{equation}
  p_N'(x)=\sum_{j=0}^N u_j\ell_j'(x).
\end{equation}
在 $x_i$ 处计算：
\begin{equation}
  p_N'(x_i)=\sum_{j=0}^N \ell_j'(x_i)u_j.
\end{equation}
定义
\begin{equation}
  \boxed{D_{ij}:=\ell_j'(x_i)}.
  \label{eq:D-definition}
\end{equation}
于是
\begin{equation}
  \boxed{p_N'(x_i)=\sum_{j=0}^N D_{ij}u_j},
  \qquad
  \boxed{\vect{u}'\approx D\vect{u}}.
\end{equation}

在 $D_{ij}$ 中：
\begin{itemize}
  \item 行指标 $i$ 表示在哪个节点 $x_i$ 计算导数；
  \item 列指标 $j$ 表示节点值 $u_j$ 对该导数的贡献；
  \item 第 $i$ 行给出 $p_N'(x_i)$ 对全部节点值的线性组合。
\end{itemize}
因此谱微分是全局的：通常 $D$ 是稠密矩阵，一个节点处的导数使用全部节点值。

\section{任意节点上的微分矩阵公式}

\subsection{节点多项式}

定义节点多项式
\begin{equation}
  \omega(x)=\prod_{k=0}^N(x-x_k).
\end{equation}
把第 $j$ 个因子提出：
\begin{equation}
  \omega(x)=(x-x_j)\prod_{k\ne j}(x-x_k).
\end{equation}
另一方面，由乘积求导法则，
\begin{equation}
  \omega'(x)=\sum_{m=0}^N\prod_{k\ne m}(x-x_k).
\end{equation}
令 $x=x_j$。当 $m\ne j$ 时，相应乘积含有 $x_j-x_j$，故整项为零；只有 $m=j$ 的一项保留：
\begin{equation}
  \omega'(x_j)=\prod_{k\ne j}(x_j-x_k).
\end{equation}
所以 Lagrange 基可改写为
\begin{equation}
  \boxed{
  \ell_j(x)=\frac{\omega(x)}{(x-x_j)\omega'(x_j)}}.
  \label{eq:lagrange-omega}
\end{equation}
这不是新定义，而是式~\eqref{eq:lagrange-basis} 的等价形式；选择它是因为求导更方便。

\subsection{非对角元素}

对式~\eqref{eq:lagrange-omega} 求导：
\begin{equation}
  \ell_j'(x)=
  \frac{\omega'(x)(x-x_j)-\omega(x)}
  {(x-x_j)^2\omega'(x_j)}.
\end{equation}
当 $i\ne j$ 且 $x=x_i$ 时，$\omega(x_i)=0$，因此
\begin{equation}
  \boxed{
  D_{ij}=\ell_j'(x_i)
  =\frac{\omega'(x_i)}{\omega'(x_j)}
  \frac{1}{x_i-x_j}},
  \qquad i\ne j.
  \label{eq:D-general}
\end{equation}
若记 $a_i:=\omega'(x_i)$，也可写成
\begin{equation}
  D_{ij}=\frac{a_i}{a_j}\frac{1}{x_i-x_j}.
\end{equation}
到这里尚未使用 Chebyshev 节点，式~\eqref{eq:D-general} 对任意互异节点都成立。

\section{Chebyshev--Gauss--Lobatto 节点}

等距节点上的高次全局插值可能在端点附近出现 Runge 振荡。Chebyshev--Gauss--Lobatto（CGL）节点取为
\begin{equation}
  \boxed{x_j=\cos\left(\frac{j\pi}{N}\right)},
  \qquad j=0,1,\ldots,N.
  \label{eq:cgl}
\end{equation}
它们在端点附近更密，并且包含
\begin{equation}
  x_0=1,\qquad x_N=-1.
\end{equation}
因此 MATLAB 数组的第一个节点是右端点，最后一个节点是左端点。这个顺序在施加非齐次边界条件时尤其重要。

\section{端点系数与交替符号的完整证明}

本节证明
\begin{equation}
  \boxed{
  \frac{\omega'(x_i)}{\omega'(x_j)}
  =(-1)^{i+j}\frac{c_i}{c_j}},
  \qquad
  c_j=\begin{cases}2,&j=0\text{ 或 }j=N,\\1,&1\le j\le N-1.\end{cases}
  \label{eq:omega-ratio}
\end{equation}

\subsection{第二类 Chebyshev 多项式}

第二类 Chebyshev 多项式满足
\begin{equation}
  U_{N-1}(\cos\theta)=\frac{\sin(N\theta)}{\sin\theta}.
  \label{eq:U-definition}
\end{equation}
对内部节点令 $\theta_j=j\pi/N$，则
\begin{equation}
  \sin(N\theta_j)=\sin(j\pi)=0,
\end{equation}
故 $U_{N-1}(x_j)=0$，$j=1,\ldots,N-1$。再乘上提供两个端点根的 $x^2-1$，定义
\begin{equation}
  q(x)=(x^2-1)U_{N-1}(x).
\end{equation}
$q$ 与 $\omega$ 都是 $N+1$ 次多项式，且具有相同的 $N+1$ 个根，所以存在非零常数 $C$ 使
\begin{equation}
  \omega(x)=Cq(x).
\end{equation}
因而导数比值中的 $C$ 会约掉：
\begin{equation}
  \frac{\omega'(x_i)}{\omega'(x_j)}
  =\frac{q'(x_i)}{q'(x_j)}.
\end{equation}

\subsection{内部节点上的导数}

求导得
\begin{equation}
  q'(x)=2xU_{N-1}(x)+(x^2-1)U_{N-1}'(x).
\end{equation}
在内部节点，第一项因 $U_{N-1}(x_j)=0$ 而消失：
\begin{equation}
  q'(x_j)=(x_j^2-1)U_{N-1}'(x_j).
\end{equation}
由 $x=\cos\theta$ 和式~\eqref{eq:U-definition}，
\begin{equation}
  U_{N-1}'(x)
  =\frac{\dfrac{\mathrm d}{\mathrm d\theta}
    \left(\dfrac{\sin(N\theta)}{\sin\theta}\right)}
   {\dfrac{\mathrm dx}{\mathrm d\theta}}.
\end{equation}
其中
\begin{equation}
  \frac{\mathrm dx}{\mathrm d\theta}=-\sin\theta
\end{equation}
且
\begin{equation}
  \frac{\mathrm d}{\mathrm d\theta}
  \left(\frac{\sin(N\theta)}{\sin\theta}\right)
  =\frac{N\cos(N\theta)\sin\theta-\sin(N\theta)\cos\theta}
  {\sin^2\theta}.
\end{equation}
在 $\theta_j=j\pi/N$ 处，利用
\begin{equation}
  \sin(N\theta_j)=0,
  \qquad \cos(N\theta_j)=(-1)^j,
\end{equation}
得到
\begin{equation}
  U_{N-1}'(x_j)=-\frac{N(-1)^j}{\sin^2\theta_j}.
\end{equation}
再使用 $x_j^2-1=-\sin^2\theta_j$，可得
\begin{equation}
  \boxed{q'(x_j)=N(-1)^j},
  \qquad 1\le j\le N-1.
  \label{eq:q-interior}
\end{equation}

\subsection{端点值为何是 $U_{N-1}(1)=N$}

$x=1$ 对应 $\theta=0$。不能把 $\theta=0$ 直接代入式~\eqref{eq:U-definition}，因为会得到 $0/0$；应取极限：
\begin{align}
  U_{N-1}(1)
  &=\lim_{\theta\to0}\frac{\sin(N\theta)}{\sin\theta}\\
  &=\lim_{\theta\to0}\frac{N\cos(N\theta)}{\cos\theta}\\
  &=N.
\end{align}
第二行使用了洛必达法则。也可以用小角近似 $\sin(N\theta)\sim N\theta$、$\sin\theta\sim\theta$ 得到同一结果。

\subsection{端点值为何是 $U_{N-1}(-1)=(-1)^{N-1}N$}

$x=-1$ 对应 $\theta=\pi$。同样必须取极限：
\begin{align}
  U_{N-1}(-1)
  &=\lim_{\theta\to\pi}\frac{\sin(N\theta)}{\sin\theta}\\
  &=\lim_{\theta\to\pi}\frac{N\cos(N\theta)}{\cos\theta}\\
  &=\frac{N\cos(N\pi)}{\cos\pi}\\
  &=\frac{N(-1)^N}{-1}\\
  &=(-1)^{N-1}N.
  \label{eq:U-minus-one}
\end{align}
最后一步成立是因为 $N+1$ 与 $N-1$ 相差 $2$，故 $(-1)^{N+1}=(-1)^{N-1}$。

还可用奇偶性检查：$U_n(-x)=(-1)^nU_n(x)$，所以
\begin{equation}
  U_{N-1}(-1)=(-1)^{N-1}U_{N-1}(1)=(-1)^{N-1}N.
\end{equation}

\subsection{端点导数与系数 2}

由于
\begin{equation}
  q'(x)=2xU_{N-1}(x)+(x^2-1)U_{N-1}'(x),
\end{equation}
在 $x=\pm1$ 处第二项含有 $x^2-1=0$，所以
\begin{align}
  q'(1)&=2U_{N-1}(1)=2N,\\
  q'(-1)&=-2U_{N-1}(-1)=2N(-1)^N.
\end{align}
与内部结果~\eqref{eq:q-interior} 合并：
\begin{equation}
  \boxed{q'(x_j)=N(-1)^jc_j},
\end{equation}
其中
\begin{equation}
  c_j=\begin{cases}2,&j=0\text{ 或 }j=N,\\1,&1\le j\le N-1.\end{cases}
\end{equation}
所以 $c_j$ 不是人为的误差修正。它只是统一记录这样一个事实：节点多项式在端点处的导数绝对值为 $2N$，在内部节点处的导数绝对值为 $N$。

最后，
\begin{align}
  \frac{\omega'(x_i)}{\omega'(x_j)}
  &=\frac{CN(-1)^ic_i}{CN(-1)^jc_j}\\
  &=(-1)^{i-j}\frac{c_i}{c_j}\\
  &=(-1)^{i+j}\frac{c_i}{c_j},
\end{align}
因为 $(-1)^{-j}=(-1)^j$，或者说 $i-j$ 与 $i+j$ 相差偶数 $2j$。这就证明了式~\eqref{eq:omega-ratio}。

\section{Chebyshev 微分矩阵}

把式~\eqref{eq:omega-ratio} 代入一般公式~\eqref{eq:D-general}，得到非对角元素
\begin{equation}
  \boxed{
  D_{ij}=\frac{c_i}{c_j}
  \frac{(-1)^{i+j}}{x_i-x_j}},
  \qquad i\ne j.
  \label{eq:D-offdiag}
\end{equation}

\subsection{对角元为何取负行和}

当 $i=j$ 时，式~\eqref{eq:D-offdiag} 的分母为零，不能直接使用。考虑常函数 $u(x)\equiv1$。其插值多项式仍为常数，导数必须为零：
\begin{equation}
  D\vect{1}=\vect{0}.
\end{equation}
这表示 $D$ 每一行元素之和为零：
\begin{equation}
  \sum_{j=0}^N D_{ij}=0.
\end{equation}
因此
\begin{equation}
  \boxed{D_{ii}=-\sum_{\substack{j=0\\j\ne i}}^N D_{ij}}.
  \label{eq:D-diagonal}
\end{equation}
负行和并不是经验技巧，而是精确保持 $\mathrm d(1)/\mathrm dx=0$ 的代数条件。

\subsection{最小例子：$N=2$}

三个节点为
\begin{equation}
  x_0=1,\qquad x_1=0,\qquad x_2=-1,
\end{equation}
微分矩阵为
\begin{equation}
  D=\begin{bmatrix}
  \frac32&-2&\frac12\\[3pt]
  \frac12&0&-\frac12\\[3pt]
  -\frac12&2&-\frac32
  \end{bmatrix}.
\end{equation}
它满足
\begin{equation}
  D\begin{bmatrix}1\\1\\1\end{bmatrix}
  =\begin{bmatrix}0\\0\\0\end{bmatrix},
\end{equation}
并且对 $u(x)=x^2$，
\begin{equation}
  D\begin{bmatrix}1\\0\\1\end{bmatrix}
  =\begin{bmatrix}2\\0\\-2\end{bmatrix},
\end{equation}
恰好等于 $u'(x)=2x$ 在节点 $1,0,-1$ 上的值。一般地，$D$ 对次数不超过 $N$ 的多项式的一阶导数是精确的（忽略浮点舍入误差）。

\section{MATLAB 实现及逐行解释}

\begin{lstlisting}[caption={Chebyshev 一阶微分矩阵}]
function [D, x] = cheb(N)
% CHEB 构造 CGL 节点和一阶谱微分矩阵。
% 若 u(j+1)=u(x_j)，则 D*u 是插值多项式在节点处的导数。

    if N == 0
        D = 0;
        x = 1;
        return
    end

    % x_0=1, x_N=-1，共 N+1 个节点
    x = cos(pi*(0:N)'/N);

    % 代码变量 c 实际存储 c_j*(-1)^j
    c = [2; ones(N-1,1); 2] .* (-1).^(0:N)';

    % dX(i,j)=x_i-x_j（数学下标比 MATLAB 下标小 1）
    X  = repmat(x,1,N+1);
    dX = X-X';

    % 构造非对角元素；eye 暂时防止对角位置除以零
    D = (c*(1./c)') ./ (dX+eye(N+1));

    % 以负行和替换临时对角元
    D = D-diag(sum(D,2));
end
\end{lstlisting}

代码中的变量 \texttt{c} 已把数学上的 $c_j$ 与 $(-1)^j$ 合并。若记
\begin{equation}
  \widetilde c_j=c_j(-1)^j,
\end{equation}
则外积 \texttt{c*(1./c)'} 的第 $(i,j)$ 个元素为
\begin{equation}
  \frac{\widetilde c_i}{\widetilde c_j}
  =\frac{c_i}{c_j}(-1)^{i-j}
  =\frac{c_i}{c_j}(-1)^{i+j}.
\end{equation}
再逐元素除以 \texttt{dX}，正好得到式~\eqref{eq:D-offdiag}。

在 \texttt{dX+eye(N+1)} 中加入单位矩阵，只是把对角分母临时从 $0$ 改为 $1$；下一行会重新计算全部对角元。若临时矩阵记为 $\widetilde D$，则
\begin{equation}
  D_{ii}=\widetilde D_{ii}-\sum_{j=0}^N\widetilde D_{ij}
  =-\sum_{j\ne i}\widetilde D_{ij},
\end{equation}
所以无论临时对角元是多少，最终都会被负行和替换。

\section{高阶导数矩阵}

由
\begin{equation}
  \vect{u}'\approx D\vect{u}
\end{equation}
再作用一次 $D$：
\begin{equation}
  \vect{u}''\approx D\vect{u}'\approx D(D\vect{u})=D^2\vect{u}.
\end{equation}
因此二阶微分矩阵写作
\begin{equation}
  \boxed{D_2=D^2}.
\end{equation}
MATLAB 中应使用矩阵乘法 \texttt{D*D} 或 \texttt{D\string^2}，不能使用 \texttt{D.\string^2}；后者只是把每个元素分别平方。

\section{经典例子：一维 Poisson 边值问题}

考虑
\begin{equation}
  u''(x)=e^{4x},\qquad -1<x<1,
  \label{eq:poisson}
\end{equation}
以及齐次 Dirichlet 边界条件
\begin{equation}
  u(-1)=u(1)=0.
\end{equation}
两次积分可得通解
\begin{equation}
  u(x)=\frac{e^{4x}}{16}+C_1x+C_2.
\end{equation}
代入两个边界条件求得
\begin{equation}
  C_1=-\frac{\sinh4}{16},
  \qquad C_2=-\frac{\cosh4}{16},
\end{equation}
故解析解为
\begin{equation}
  \boxed{u(x)=\frac{e^{4x}-x\sinh4-\cosh4}{16}}.
\end{equation}

\subsection{离散方程与边界条件}

在全部节点上，式~\eqref{eq:poisson} 离散为
\begin{equation}
  D^2\vect{u}\approx\vect{f}.
\end{equation}
记内部未知量为
\begin{equation}
  \vect{u}_{\mathrm{in}}=
  \begin{bmatrix}u_1&\cdots&u_{N-1}\end{bmatrix}^{\mathsf T}.
\end{equation}
内部行可以写成
\begin{equation}
  D^2_{\mathrm{in},\mathrm{in}}\vect{u}_{\mathrm{in}}
  +D^2_{\mathrm{in},0}u_0
  +D^2_{\mathrm{in},N}u_N
  =\vect{f}_{\mathrm{in}}.
\end{equation}
本例中 $u_0=u(1)=0$、$u_N=u(-1)=0$，所以边界列对应的项消失，只需解
\begin{equation}
  D^2_{\mathrm{in},\mathrm{in}}\vect{u}_{\mathrm{in}}
  =\vect{f}_{\mathrm{in}}.
\end{equation}

\begin{lstlisting}[caption={Poisson 边值问题求解程序}]
clear; clc; close all;

N = 16;
[D,x] = cheb(N);
D2 = D^2;

% MATLAB 的 2:N 对应数学下标 1,...,N-1
A = D2(2:N,2:N);
f = exp(4*x(2:N));

u_inner = A\f;
u = [0;u_inner;0];

u_exact = (exp(4*x)-x*sinh(4)-cosh(4))/16;
err_inf = norm(u-u_exact,inf);
fprintf('N=%d, 最大节点误差=%.3e\n',N,err_inf);

% x 原本从 1 排到 -1；绘图前升序排列
[x_plot,idx] = sort(x);
plot(x_plot,u_exact(idx),'k-','LineWidth',1.5); hold on;
plot(x_plot,u(idx),'ro','MarkerSize',6);
grid on;
xlabel('x'); ylabel('u(x)');
legend('解析解','谱方法节点解','Location','best');
\end{lstlisting}

\subsection{非齐次边界条件}

若
\begin{equation}
  u(1)=\alpha,\qquad u(-1)=\beta,
\end{equation}
由于节点顺序为 $x_0=1$、$x_N=-1$，内部右端项必须改为
\begin{equation}
  \vect{f}_{\mathrm{in}}
  -D^2_{\mathrm{in},0}\alpha
  -D^2_{\mathrm{in},N}\beta.
\end{equation}
对应代码为
\begin{lstlisting}
rhs = f - D2(2:N,1)*alpha - D2(2:N,N+1)*beta;
u_inner = D2(2:N,2:N)\rhs;
u = [alpha;u_inner;beta];
\end{lstlisting}
非齐次边界下不能只删掉首尾两列而不修改右端项。

\section{验证微分矩阵}

在求解微分方程前，先验证 $D$ 的基本性质。对于次数不超过 $N$ 的多项式，$D$ 应在浮点舍入范围内给出精确导数。

\begin{lstlisting}[caption={微分矩阵自检}]
clear; clc;
N = 8;
[D,x] = cheb(N);

err_const = norm(D*ones(size(x)),inf);
err_x     = norm(D*x-ones(size(x)),inf);
err_x2    = norm(D*(x.^2)-2*x,inf);
err_xN    = norm(D*(x.^N)-N*x.^(N-1),inf);

fprintf('常函数误差 = %.3e\n',err_const);
fprintf('x 导数误差 = %.3e\n',err_x);
fprintf('x^2 导数误差 = %.3e\n',err_x2);
fprintf('x^N 导数误差 = %.3e\n',err_xN);
\end{lstlisting}

还可检查
\begin{equation}
  \lVert D\vect{1}\rVert_\infty
\end{equation}
是否接近机器精度。若它明显不小，通常意味着差矩阵方向、权重比值或对角元处理有误。

\section{谱收敛实验}

对光滑或解析函数，Chebyshev 谱方法的误差通常随 $N$ 快速下降，常表现为指数型或几何型收敛，直到舍入误差和矩阵条件数开始主导。

\begin{lstlisting}[caption={误差随 $N$ 的变化}]
Nvec = 2:2:40;
err = zeros(size(Nvec));

for k = 1:numel(Nvec)
    N = Nvec(k);
    [D,x] = cheb(N);
    D2 = D^2;
    f = exp(4*x(2:N));
    u = [0;D2(2:N,2:N)\f;0];
    ue = (exp(4*x)-x*sinh(4)-cosh(4))/16;
    err(k) = norm(u-ue,inf);
end

semilogy(Nvec,err,'o-');
grid on;
xlabel('N'); ylabel('最大节点误差');
title('Chebyshev 谱配置法的收敛');
\end{lstlisting}

半对数坐标上若误差在一段区间内近似直线下降，说明误差大致满足 $C\rho^{-N}$，其中 $\rho>1$。误差触及约 $10^{-14}$ 附近后不再下降，并不表示理论失效，而通常是双精度舍入误差已经成为主要误差源。

\section{常见混淆与排错}

\begin{longtable}{>{\raggedright\arraybackslash}p{0.27\textwidth}p{0.65\textwidth}}
\toprule
混淆 & 正确理解\\
\midrule
\endhead
$\ell_j(x_i)=\delta_{ij}$ 是否说明 $\ell_j$ 恒为零？
& 不说明。它只规定有限个节点上的取值；在自己的节点处为 $1$，节点之间一般非零。\\

$D$ 是原函数的精确导数算子吗？
& $D$ 精确求的是插值多项式 $p_N$ 在节点上的导数。它对原函数导数的准确程度取决于插值误差和函数光滑性。\\

$c_0=c_N=2$ 是经验修正吗？
& 不是。它来自 $q'(\pm1)$ 的绝对值为 $2N$，而内部 $q'(x_j)$ 的绝对值为 $N$。\\

为什么代码里的 \texttt{c} 有正负号？
& 代码把数学上的端点系数 $c_j$ 和交替因子 $(-1)^j$ 合并为一个向量。\\

为什么 $D$ 是稠密的？
& 因为每个 Lagrange 基函数是全局多项式，一个节点值会影响整条插值曲线。\\

为什么 $D^2$ 不是逐元素平方？
& 连续求导对应线性算子的复合，因此是矩阵乘法 $D D$。\\

为什么结果符号全部相反？
& 最常见原因是把 $x_i-x_j$ 写成了 $x_j-x_i$，或节点顺序与公式约定不一致。用 $D\vect{x}=\vect{1}$ 检查。\\

为什么非齐次边界结果不对？
& 删除边界列前，必须把这些列乘已知边界值后移到右端。\\
\bottomrule
\end{longtable}

\section{符号表}

\begin{longtable}{>{\centering\arraybackslash}p{0.18\textwidth}p{0.72\textwidth}}
\toprule
符号 & 含义\\
\midrule
\endhead
$N$ & 插值多项式最高次数；节点总数为 $N+1$。\\
$i,j,k,m$ & 节点或连乘、求和指标，数学下标从 $0$ 开始。\\
$x_j$ & 第 $j$ 个节点；CGL 节点为 $\cos(j\pi/N)$。\\
$u_j$ & $u(x_j)$，原函数在第 $j$ 个节点上的值。\\
$p_N(x)$ & 穿过全部节点值的次数不超过 $N$ 的插值多项式。\\
$\ell_j(x)$ & 第 $j$ 个 Lagrange 基函数。\\
$\delta_{ij}$ & Kronecker 符号；$i=j$ 时为 $1$，否则为 $0$。\\
$\omega(x)$ & 节点多项式 $\prod_{k=0}^N(x-x_k)$。\\
$D_{ij}$ & $\ell_j'(x_i)$，一阶微分矩阵第 $i$ 行第 $j$ 列。\\
$U_n(x)$ & 第二类 Chebyshev 多项式。\\
$c_j$ & 端点修正系数；端点为 $2$，内部为 $1$。\\
$D^2$ & 一阶微分矩阵的矩阵平方，用于二阶导数。\\
$\lVert\cdot\rVert_\infty$ & 向量分量绝对值的最大值。\\
\bottomrule
\end{longtable}

\section{建议练习}

\begin{enumerate}[label=\arabic*.]
  \item 手算 $N=2$ 时全部 $D_{ij}$，并验证 $D\vect{1}=0$、$D\vect{x}=\vect{1}$。
  \item 用程序验证 $D(x^m)=m x^{m-1}$，其中 $m=0,1,\ldots,N$。
  \item 把方程改为 $-u''+u=f$，写出离散矩阵 $-D^2+I$。
  \item 取非齐次边界 $u(1)=1$、$u(-1)=2$，验证右端修正项。
  \item 比较 Chebyshev 谱方法与二阶中心差分的误差随节点数变化。
  \item 解释当函数不光滑时，为什么谱收敛速度会显著下降。
\end{enumerate}

\section{核心结论}

整套方法最重要的三条公式是
\begin{equation}
  \boxed{D_{ij}=\ell_j'(x_i)},
\end{equation}
它定义了微分矩阵的本质；
\begin{equation}
  \boxed{
  D_{ij}=\frac{c_i}{c_j}
  \frac{(-1)^{i+j}}{x_i-x_j},\quad i\ne j},
\end{equation}
它给出了 CGL 节点上的非对角元素；以及
\begin{equation}
  \boxed{D_{ii}=-\sum_{j\ne i}D_{ij}},
\end{equation}
它由常函数导数为零确定对角元素。

从思想上说，谱配置法完成了如下转换：
\begin{equation*}
\boxed{\text{微积分中的求导}\quad\longrightarrow\quad\text{线性代数中的矩阵乘法}}.
\end{equation*}

\end{document}
